\chapter{SCIENTIFIC LITERATURE REVIEW}

\section{Gas While Drilling (GWD) and Mud Logging Fundamentals}

In petroleum drilling engineering, mud logging represents the continuous surface extraction, monitoring, and geochemical analysis of rock cuttings and formation fluids carried out of the borehole by the circulating drilling fluid \citep{ref9, ref12}. As the rotary drill bit grinds through subsurface geological strata, hydrocarbons trapped within the rock pore spaces are crushed and liberated into the drilling mud. 

At the surface, a mechanical gas trap agitator installed at the "possum belly" of the shale shaker extracts the volatile gas mixture from the mud stream. A continuous pneumatic sample line transports the gas to a field logging unit equipped with a high-speed gas chromatograph (GC) and a total hydrocarbon analyzer (flame ionization detector, FID) \citep{ref12, haworth1985}. The chromatographic analysis decomposes the liberated hydrocarbon gas into its individual alkane constituents:
\begin{itemize}
    \item \textbf{Methane ($C_1$, $\text{CH}_4$):} The lightest, most volatile alkane molecule, characteristic of non-associated dry natural gas reservoirs.
    \item \textbf{Ethane ($C_2$, $\text{C}_2\text{H}_6$):} Light hydrocarbon gas fraction, associated with gas-condensate and volatile oil caps.
    \item \textbf{Propane ($C_3$, $\text{C}_3\text{H}_8$):} Intermediate-weight alkane, serving as a primary discriminator between gaseous and liquid hydrocarbon phases.
    \item \textbf{Iso-butane ($iC_4$) and Normal-butane ($nC_4$, $\text{C}_4\text{H}_{10}$):} Medium-weight alkanes indicating liquid-associated hydrocarbons and crude oil columns.
    \item \textbf{Iso-pentane ($iC_5$) and Normal-pentane ($nC_5$, $\text{C}_5\text{H}_{12}$):} Heavy hydrocarbon gas fractions, strongly indicative of liquid crude oil accumulations.
    \item \textbf{Total Gas ($TG$):} The aggregate sum of all hydrocarbon gases liberated into the drilling mud, representing the overall magnitude of the formation gas show \citep{ref14}.
\end{itemize}

Because GWD chromatographic data is collected in real time during the initial drilling pass, analyzing the relative proportions, ratios, and compositional trends of these alkane fractions provides the earliest diagnostic insights into reservoir fluid phase, fluid contacts, and formation productivity \citep{haworth1985, ref12}.

\section{Critical Review of Existing Formation Evaluation Software}

The interpretation of well log and mud log data has historically relied on the expertise of petrophysicists supported by specialized computer-aided formation evaluation software \citep{ref8, wood2020}. Over past decades, commercial software suites and programmatic libraries have transformed wellsite workflows by replacing paper strip logs with digital visualization environments. However, an objective examination reveals critical limitations:

\subsection{Commercial Workstation Suites}
\begin{enumerate}
    \item \textbf{Schlumberger (SLB) Techlog and AspenTech Geolog:}
    \begin{itemize}
        \item \textit{Capabilities:} High-end enterprise platforms providing advanced multi-track log visualization, borehole image processing, geomechanical modeling, and integration with large-scale corporate databases \citep{wood2020}.
        \item \textit{Limitations:} Heavyweight workstation architecture requiring complex installations, high-performance hardware, and multi-thousand-dollar annual licensing fees (frequently exceeding \$20,000 to \$50,000 per seat) \citep{wood2020}. Crucially, with respect to GWD mud log data, they function primarily as passive visualizers. Calculating multi-ratio petrophysical indicators or executing automated fluid predictions requires writing complex routines in specialized proprietary scripting consoles (e.g., Techlog Python Script Editor or Geolog Loglan macros) or compiling C++ DLL plugins.
    \end{itemize}
    
    \item \textbf{Advanced Logic Technology WellCAD and RockWare LogPlot:}
    \begin{itemize}
        \item \textit{Capabilities:} Industry-standard geotechnical drafting engines providing flexible visual layouts for lithological patterns, core descriptions, and standardized wellsite composite logs.
        \item \textit{Limitations:} Strictly visual plotting engines with zero built-in petrophysical calculation logic. They plot raw gas values ($C_1 - C_5$) as static visual lines but possess no automated algorithms to derive ratio matrices or forecast hydrocarbon fluid regimes.
    \end{itemize}
\end{enumerate}

\subsection{Open-Source and Programmatic Log Plotters}
In the open-source software ecosystem, several Python libraries (such as \texttt{lasio}, \texttt{welly}, and standard \texttt{matplotlib/plotly} scripts) have been developed for petrophysical data manipulation \citep{wood2020}:
\begin{itemize}
    \item \textit{Capabilities:} Free, scriptable parsing of Log ASCII Standard (\texttt{.las}) files and programmatic 2D line plotting without commercial licensing costs.
    \item \textit{Limitations:} They lack an integrated graphical user interface (GUI), requiring users to write manual Python code for every file parsing, ratio calculation, and visual track configuration. Furthermore, they contain no automated GWD multi-ratio interpretation pipelines or deterministic decision-tree logic, serving only as low-level scripting building blocks.
\end{itemize}

\section{Theoretical and Mathematical Formulations of Petrophysical Ratios}

To extract diagnostic intelligence from raw chromatographic data, petroleum engineers and petrophysicists have established deterministic ratio formulations. The calculation engine developed in this research implements 16 derived petrophysical indicators governed by canonical literature:

\subsection{Classic and Expanded Pixler Ratios \citep{ref13}}
The Pixler method evaluates light-to-medium hydrocarbon fractions to delineate gas from oil regimes:
\begin{equation}
    R_1 = \frac{C_1}{C_2}, \quad R_2 = \frac{C_1}{C_3}, \quad R_3 = \frac{C_1}{C_4}, \quad R_4 = \frac{C_1}{C_5}
\end{equation}
In accordance with \citet{ref13}, the primary diagnostic cutoff boundaries for $R_1$ and crossplot slopes are defined as:
\begin{itemize}
    \item \textbf{Non-Associated Dry Gas:} $R_1 > 65.0$.
    \item \textbf{Gas / Gas-Condensate Pay:} $15.0 \le R_1 \le 65.0$.
    \item \textbf{Oil Pay Zone (Medium to Light Crude):} $2.0 \le R_1 < 15.0$.
    \item \textbf{Water / Non-Productive Formation:} $R_1 < 2.0$.
    \item \textbf{Pixler Slope Rule ($\Delta R$):} A progressive positive slope ($R_1 < R_2 < R_3 < R_4$) confirms commercial productive pay. A flat, irregular, or negative slope indicates tight, low-permeability rock or water-bearing intervals.
\end{itemize}

\subsection{Expanded Butane Isomer Multipliers}
Tracks methane against separated butane isomers to capture sensitivity to heavier, liquid-associated hydrocarbon fractions:
\begin{equation}
    \text{Ratio}_{\text{iC4}} = \frac{C_1}{iC_4}, \quad \text{Ratio}_{\text{nC4}} = \frac{C_1}{nC_4}
\end{equation}
\begin{itemize}
    \item \textbf{Gas Pay Threshold:} $\text{Ratio}_{\text{nC4}} \ge 30.0$ (methane dominates over butane isomers).
    \item \textbf{Oil Pay Threshold:} $\text{Ratio}_{\text{nC4}} < 30.0$ (enriched heavy butane indicates liquid crude oil).
\end{itemize}

\subsection{Total Gas Volume ($TG$) and Gas Dryness Ratio ($DR$) \citep{ref14}}
Calculated by summing all recorded hydrocarbon gas fractions when Total Gas is not pre-recorded:
\begin{equation}
    TG = C_1 + C_2 + C_3 + iC_4 + nC_4 + iC_5 + nC_5
\end{equation}
The Dryness Ratio ($DR$) quantifies the methane fraction within the total hydrocarbon gas mixture:
\begin{equation}
    DR = \frac{C_1}{TG}
\end{equation}
\begin{itemize}
    \item \textbf{Dry Gas Regime:} $DR \ge 0.85$ (methane concentration $\ge 85\%$).
    \item \textbf{Wet Gas / Light Volatile Oil:} $0.50 \le DR < 0.85$.
    \item \textbf{Heavy Crude Oil / Liquid-Rich Zone:} $DR < 0.50$.
    \item \textbf{Baseline Show Threshold:} $TG \ge 300\text{ ppm}$ or $C_1 \ge 200\text{ ppm}$ (intervals below threshold $\rightarrow$ \textbf{No Show / Barren Rock}).
\end{itemize}

\subsection{Carbon-Weighted Density Index ($I_{\text{carbon}}$)}
Quantifies carbon density across the gas stream, scaled by carbon atom count:
\begin{equation}
    I_{\text{carbon}} = \frac{TG}{C_1 + 2C_2 + 3C_3 + 4iC_4 + 4nC_4 + 5iC_5 + 5nC_5}
\end{equation}
\begin{itemize}
    \item \textbf{Light Hydrocarbon Gas:} $I_{\text{carbon}} \ge 0.80$.
    \item \textbf{Heavy Liquid Column:} $I_{\text{carbon}} < 0.80$.
\end{itemize}

\subsection{Expanded Haworth Show Evaluation Ratios \citep{haworth1985, ref12}}
These three expanded parameters form the foundational multi-ratio fluid typing framework (Wetness $W_h$, Balance $B_h$, and Character $C_h$):
\begin{equation}
    W_h = \left( \frac{C_2 + C_3 + iC_4 + nC_4 + iC_5 + nC_5}{TG} \right) \times 100\%
\end{equation}
\begin{equation}
    B_h = \frac{C_1 + C_2}{C_3 + iC_4 + nC_4 + iC_5 + nC_5}
\end{equation}
\begin{equation}
    C_h = \frac{iC_4 + nC_4 + iC_5 + nC_5}{C_3}
\end{equation}
According to canonical Haworth et al. criteria:
\begin{itemize}
    \item \textbf{Haworth Wetness Ratio ($W_h$):}
    \begin{itemize}
        \item $W_h < 0.5\%$: Very dry gas, non-commercial tight formation, or barren shale.
        \item $0.5\% \le W_h < 17.5\%$: Gas Pay regime ($0.5\% \le W_h < 5.0\%$ indicates dry gas; $5.0\% \le W_h < 17.5\%$ indicates wet gas/condensate).
        \item $17.5\% \le W_h < 40.0\%$: Oil Pay regime (light to medium crude oil).
        \item $W_h \ge 40.0\%$: Heavy, low-gravity residual or dead oil.
    \end{itemize}
    \item \textbf{Haworth Balance Ratio ($B_h$):}
    \begin{itemize}
        \item $B_h > 100.0$: Very dry gas.
        \item $15.0 \le B_h \le 100.0$: Gas pay column.
        \item $0.5 \le B_h < 15.0$: Oil pay column.
        \item $B_h < 0.5$: Heavy residual oil, water, or barren rock.
    \end{itemize}
    \item \textbf{Haworth Character Ratio ($C_h$):}
    \begin{itemize}
        \item $C_h < 0.5$: Gas phase (light alkane dominance).
        \item $C_h \ge 0.5$: Oil phase (heavier butane/pentane enrichment).
    \end{itemize}
\end{itemize}

\subsection{Composite Screening Indicators ($GOW, WBS, GOR$) \citep{whittaker1991, spe2020}}
Continuous composite indicators isolate heavy fluid concentrations and decouple total gas fluctuations:
\begin{equation}
    \text{GOW} = \frac{(C_3 + iC_4 + nC_4 + iC_5 + nC_5) \cdot TG}{TG}, \quad \text{GOW}_{\text{noTG}} = \frac{C_1 \cdot C_2}{C_3 \cdot C_4 \cdot C_5}
\end{equation}
\begin{equation}
    WBS = \frac{\log_{10}(B_h) - 0.903}{2.097} - \frac{\log_{10}(W_h)}{2}
\end{equation}
\begin{equation}
    \text{SPE } GOR = \text{Empirical Gas-to-Oil Ratio Proxy}
\end{equation}
\begin{itemize}
    \item \textbf{Whittaker Score ($WBS$):} $WBS > 2.0 \rightarrow$ Gas Pay; $0.0 \le WBS \le 2.0 \rightarrow$ Oil Pay; $WBS < 0.0 \rightarrow$ Water / Non-Productive.
    \item \textbf{Composite $GOW_{\text{noTG}}$:} $GOW > 10.0 \rightarrow$ Gas Pay; $1.0 \le GOW \le 10.0 \rightarrow$ Oil Pay; $GOW < 1.0 \rightarrow$ Water / Residual.
    \item \textbf{SPE GOR Proxy:} $GOR \ge 3000\text{ scf/bbl} \rightarrow$ Gas / Condensate; $GOR < 3000\text{ scf/bbl} \rightarrow$ Oil column.
\end{itemize}

Table~\ref{tab:lit_thresholds_chap3} synthesizes the complete matrix of canonical literature threshold criteria establishing the baseline decision rules implemented in the engine.

\begin{table}[H]
\centering
\caption{Canonical Literature Petrophysical Thresholds and Baseline Decision Rules}
\label{tab:lit_thresholds_chap3}
\renewcommand{\arraystretch}{1.3}
\resizebox{\textwidth}{!}{%
\begin{tabular}{|l|l|c|c|c|l|}
\hline
\textbf{Indicator} & \textbf{Mathematical Formulation} & \textbf{Gas Pay Range} & \textbf{Oil Pay Range} & \textbf{Water / Non-Pay} & \textbf{Primary Literature Source} \\
\hline
Pixler Ratio 1 ($R_1$) & $C_1 / C_2$ & $15.0 \le R_1 \le 65.0$ & $2.0 \le R_1 < 15.0$ & $R_1 < 2.0$ (or $>65$) & \citet{ref13} \\
Pixler Slope ($\Delta R$) & Progression $R_1 < R_2 < R_3 < R_4$ & Positive Slope & Positive Slope & Flat / Negative & \citet{ref13} \\
Haworth Wetness ($W_h$) & $\frac{\sum_{k=2}^5 C_k}{\sum_{k=1}^5 C_k} \times 100\%$ & $0.5\% \le W_h < 17.5\%$ & $17.5\% \le W_h < 40.0\%$ & $W_h < 0.5\%$ / $\ge 40\%$ & \citet{haworth1985} \\
Haworth Balance ($B_h$) & $\frac{C_1 + C_2}{C_3 + iC_4 + nC_4 + C_5}$ & $15.0 \le B_h \le 100.0$ & $0.5 \le B_h < 15.0$ & $B_h < 0.5$ (or $>100$) & \citet{haworth1985} \\
Haworth Character ($C_h$) & $\frac{iC_4 + nC_4 + iC_5 + nC_5}{C_3}$ & $C_h < 0.5$ & $C_h \ge 0.5$ & Non-diagnostic & \citet{haworth1985} \\
Dryness Ratio ($DR$) & $C_1 / TG$ & $DR \ge 0.85$ & $0.50 \le DR < 0.85$ & $DR < 0.50$ & \citet{ref14} \\
Whittaker Score ($WBS$) & Standardized Log Ratio Score & $WBS > 2.0$ & $0.0 \le WBS \le 2.0$ & $WBS < 0.0$ & \citet{whittaker1991} \\
Composite $GOW_{\text{noTG}}$ & $(C_1 \cdot C_2) / (C_3 \cdot C_4 \cdot C_5)$ & $GOW > 10.0$ & $1.0 \le GOW \le 10.0$ & $GOW < 1.0$ & \citet{whittaker1991} \\
SPE GOR Proxy & Empirical $GOR$ Formula & $GOR \ge 3000\text{ scf/bbl}$ & $GOR < 3000\text{ scf/bbl}$ & Non-commercial & \citet{spe2020} \\
\hline
\multicolumn{6}{|l|}{\textbf{Baseline Gas Show Detection Gate:} $TG \ge 300\text{ ppm}$ or $C_1 \ge 200\text{ ppm}$ (Intervals below threshold $\rightarrow$ \textbf{No Show / Barren Rock})} \\
\hline
\end{tabular}%
}
\end{table}

\section{Machine Learning Approaches vs. Regulatory PRMS Audit Constraints}

Parallel to classical formation evaluation techniques, the application of machine learning (ML) and deep learning (DL) to subsurface prediction has grown substantially since its early introduction to petroleum engineering \citep{ref1}. The conceptual groundwork for applying artificial intelligence techniques to oilfield problems was laid by \citet{ref1}, who demonstrated that Artificial Neural Networks (ANN) could approximate complex, non-linear mappings between well log inputs and reservoir properties that resist explicit analytical formulation. 

Over the subsequent decades, data-driven methods have expanded into rock mechanical parameter prediction \citep{ref5}, automated curve reconstruction and missing log completion using gradient-boosted decision trees \citep{ref4}, and flexible graph neural networks (GNN) such as FlexLogNet \citep{ref6}. While these statistical models achieve high empirical curve-fitting performance, they present fundamental limitations when applied to real-time formation evaluation:
\begin{enumerate}
    \item \textbf{Opacities and Black-Box Nature:} Neural networks and ensemble trees distribute decision logic across millions of statistical weights and non-linear activation functions, obscuring physical causality.
    \item \textbf{SPE PRMS Chapter 4 Reserve Certification Mandates:} Under international reserve evaluation guidelines established by the Society of Petroleum Engineers (SPE), World Petroleum Council (WPC), American Association of Petroleum Geologists (AAPG), Society of Petroleum Evaluation Engineers (SPEE), and Society of Exploration Geophysicists (SEG), all procedures determining commercial hydrocarbon reserves and net pay intervals must be 100\% mathematically traceable, deterministic, and auditable \citep{ref2, ref3}. Black-box machine learning predictions are legally disqualified during formal commercial audits and public disclosures.
    \item \textbf{Training Data Scarcity in Wildcat Wells:} Supervised ML models require massive historical training datasets with verified core or production tests, which are unavailable when drilling exploratory wildcat wells in frontier basins.
\end{enumerate}

These regulatory constraints justify the development of a transparent, deterministic rule-based calculation engine that delivers automated multi-ratio intelligence while strictly preserving PRMS auditability.

\section{Software Engineering Framework and Core Technologies}

The software engineering methodology follows an adapted Software Development Life Cycle (SDLC) tailored for scientific computing and operational decision support systems:
\begin{enumerate}
    \item \textbf{Python 3 Computing Ecosystem:} The core computational engine is implemented in Python, leveraging \textbf{NumPy} for vectorized array processing and contiguous memory operations, and \textbf{Pandas} for high-performance tabular data frame manipulation.
    \item \textbf{Dash Reactive Web Framework:} Built by Plotly, Dash integrates Flask, React.js, and Plotly.js into a reactive Python environment. Interactive callbacks (\texttt{@app.callback}) execute asynchronous state updates in memory without requiring full-page reloads.
    \item \textbf{Plotly WebGL Visualization Engine:} Plotly provides hardware-accelerated WebGL rendering for continuous vertical subplots, generating synchronized depth axes, hover crosshairs, and dynamic shaded area fills without browser DOM degradation.
    \item \textbf{Vercel Serverless Cloud Deployment:} To provide global, zero-installation access, the application is packaged as a serverless application deployed on the \textbf{Vercel} edge platform, leveraging edge network caching for sub-second delivery over high-latency field connections.
\end{enumerate}